波算子在 $(t,x,s_0)$ 中为 $2\partial_t\partial_{s_0}+\Delta_x$，因此在频率 $\kappa$（$s_0$ 变量中）处变为 $2\mathrm i\kappa\partial_t+\Delta_x$。把 $T=(t-s_0)/\sqrt2$ 与 $b=(t+s_0)/\sqrt2$ 代入相位，因 $|\xi|^2-\xi_b^2=|\xi_x|^2$ 而得所述模式。$\mathcal P(\partial_{s_0}\psi,\nabla_x\psi)^{\mathsf t}$ 的第一分量为 $(2\partial_t\partial_{s_0}+\Delta_x)\psi$，其他分量因导数可交换而抵消。直接计算 $C(p)$ 得到所列分块矩阵。它的二次型为所述平方，核由 $E(v)$ 的倍数组成。作用在该模式上时，$\mathcal P$ 等同于乘以 $\mathrm i C(p)$，由此得到向量解。

在物理坐标中按 $-\mathrm i\rho(1,\omega)$ 平移，其中 $\rho>0$ 和 $|\omega|=1$，等同于将模式乘以 $e^{\rho(\xi\cdot\omega-|\xi|)}=e^{-\rho|\xi||\xi/|\xi|-\omega|^2/2}$。在 $|\xi|\asymp N$ 和 $\rho N\gg1$ 处，其角宽量级为 $(\rho N)^{-1/2}$。二次通量 $q(w)^\mu=w^{\mathsf t}C^\mu w/2$ 在零坐标标架中为 $(w_0^2,w_0w',-|w'|^2/2)$。在物理坐标中，其时间分量与空间部分的长度均等于 $(w_0^2+|w'|^2/2)/\sqrt2$。在 $w=E(v)$ 处，物理通量为 $(|\xi|,\xi)/\kappa$。

在两个通量之和中，由对称性可知 $z^*C^\mu z$ 的混合项抵消。对每个实状态使用 $-C(p)$ 的二次型，得到 $p\cdot q=-\kappa|\zeta-vR|^2/2$。将 $\kappa^d e^{-\kappa|\zeta-vR|^2}$ 的积分在 $2d$ 个实坐标中平移并伸缩，得到 $\pi^d$。
